% --- LaTeX Essay Template - S. Venkatraman ---

% --- Set document class and font size ---

\documentclass[letterpaper, 11pt]{article}

% --- Package imports ---

% lipsum is just for generating placeholder text and can be removed
\usepackage{hyperref, lipsum} 

% --- Fonts (Set to Palatino or Times New Roman if desired) ---

% \usepackage{newpxtext, newpxmath}
% \usepackage{times}

% --- Page layout settiings ---

% Set page margins
\usepackage[left=1.2in, right=1.2in, top=.9in, bottom=.9in]{geometry}

% Anchor footnotes to the bottom of the page
\usepackage[bottom]{footmisc}
\usepackage{setspace}
\setstretch{1.15}
% Set line spacing
\renewcommand{\baselinestretch}{1.2}

% Use this for a one-off '*' footnote (e.g., a URL in a heading)
\newcommand{\starfootnote}[1]{%
\begingroup
\renewcommand{\thefootnote}{\fnsymbol{footnote}}%
\footnote{#1}%
\endgroup
\addtocounter{footnote}{-1}%
}

% Set spacing between paragraphs
\setlength{\parskip}{2mm}

% --- Page formatting settings ---



% --- Essay title and author ---
\title{\vspace{-1.5cm} Summary of \textit{Veridical Data Science}, Chapters 13 \footnote{\url{https://vdsbook.com/}}}

\date{\today}

% --- Main text ---
% --- Document starts here ---

\begin{document}
% Set link colors for labeled items and URLs (blue) 
\hypersetup{colorlinks=true, linkcolor=blue, urlcolor=blue}
\maketitle

\section*{Chapter 13: Decision Trees and Random Forests}
Chapter 13 introduces the uncertainty in the data science process that arises from the many choices and judgments made along the way, such as in data cleaning, preprocessing, and algorithm selection.
The traditional machine-learning approach, which selects the best algorithm trained on one cleaned and preprocessed dataset, which cannot capture this uncertainty, and thus can lead to overconfident predictions.
The chapter then presents the three different frameworks, which provides a structured approach to making these choices and judgments in a way that leads to more reliable and trustworthy predictions.

The first is the single PCS fit, which selects the best-performing predictive score on the validation set.
The second is the PCS ensemble, which aggregates predictions from all fits that pass problem-specific predictability screening, such as by averaging for continuous responses or majority voting for binary responses.
However, these two approaches cannot deliver the level of uncertainty quantification that is often desired in practice, because they only give a single value.
The third is the prediction perturbation intervals (PPIs), which is designed for continuous responses and communicates uncertainty by reporting an interval rather than a single value.
The chapter emphasizes that PPIs must be calibrated to achieve a desired coverage level, because raw prediction quantiles alone do not guarantee correct coverage.
The chapter also stresses that the test set should be reserved for one final independent evaluation after the final prediction approach has been chosen, and should not be used for model selection or tuning, as this can lead to overfitting and overly optimistic performance estimates.

% --- Document ends here ---
\end{document}
